% ==================================================================
% id:       app-018
% topic:    שימושי הנגזרת
% subtopic: פולינום טיילור
% solution: full
% ==================================================================
\bE
 \item ידוע שהגבול $\ds{\lim_{x\to 0}\frac{e^{2x}-\cos x-\sin(2x)+ax^2}{x^3}}$ קיים וסופי. מצאו את $a$  ואת הגבול.
\hint{פתחו את כל הפונקציות עד סדר $3$ ודרשו שהמקדמים של החזקות הנמוכות מ-$x^3$ במונה יתאפסו.}
\begin{solbox}
נציג
$$,e^{2x}=1+2x+2x^2+\frac{4}{3}x^3+o(x^3),\ \cos x=1-\frac{x^2}{2}+o(x^3)$$
$$.\sin(2x)=2x-\frac{4}{3}x^3+o(x^3)$$
לכן
\[
e^{2x}-\cos x-\sin(2x)=\frac{5}{2}x^2+\frac{8}{3}x^3+o(x^3).
\]
מכאן
\begin{align*}
    \frac{e^{2x}-\cos x-\sin(2x)+ax^2}{x^3}&=\ \frac{\left(\frac{5}{2}+a\right)x^2+\frac{8}{3}x^3+o(x^4)}{x^3}.
\end{align*}
לקיום הגבול נדרוש כי
$$,a=-\frac{5}{2}$$
ואז
\begin{align*}
    \lim_{x\to 0}\frac{e^{2x}-\cos x-\sin(2x)+ax^2}{x^3}&=\ \lim_{x\to 0}\frac{\frac{8}{3}x^3+o(x^4)}{x^3}
    =\ \frac{8}{3}.
\end{align*}
\end{solbox}
 \item ידוע שהגבול $\ds{\lim_{x \to 0}\frac{x-\lr{a+b\cos x}\sin x}{\sin^5x}}$ קיים. מצאו את $a$, את $b$ ואת הגבול.
\hint{פתחו את המונה עד סדר $5$, והשתמשו בכך ש-$\frac{\sin x}{x}\to1$ כדי להחליף את המכנה ב-$x^5$. דרשו שהמקדמים של החזקות הנמוכות יתאפסו.}
\begin{solbox}
נציג
\[
,\sin x=x-\frac{x^3}{6}+\frac{x^5}{120}+o(x^5),\
\cos x=1-\frac{x^2}{2}+\frac{x^4}{24}+o(x^4)
\]
מכאן
\[
(a+b\cos x)\sin x=(a+b)x-\frac{a+4b}{6}x^3+\frac{a+16b}{120}x^5+o(x^5).
\]
ולכן
\[
x-(a+b\cos x)\sin x=\bigl(1-(a+b)\bigr)x+\frac{a+4b}{6}x^3-\frac{a+16b}{120}x^5+o(x^5).
\]
נסמן
$$A=1-(a+b),\ B=\frac{a+4b}{6},\ C=\frac{a+16b}{120}$$
מתקיים
\begin{align*}
    \lim_{x \to 0}\frac{x-\lr{a+b\cos x}\sin x}{\sin^5x}&=\ \lim_{x \to 0}\frac{Ax+Bx^3-Cx^5+o(x^5)}{\sin^5x}\\
    &=\ \lim_{x \to 0}\frac{Ax+Bx^3-Cx^5+o(x^5)}{\color{red}x^5}{\color{red}\cdot\frac{x^5}{\sin^5x}}.
\end{align*}
ע"מ שהגבול יהיה סופי נדרוש כי
$$,A=0, B=0$$
כלומר
$$,a=\frac{4}{3},\ b=-\frac{1}{3}$$
ואז
$$.\lim_{x \to 0}\frac{x-\lr{a+b\cos x}\sin x}{\sin^5x}=-C=\frac{1}{30}$$
\end{solbox}
 \item ידוע שהגבול $\ds{\lim_{x\to 0}\frac{e^{2x}\sin x+ax+bx^2}{x^3}}$ קיים וסופי. מצאו את $a$, $b$ ואת הגבול.
\begin{solbox}
בדומה לתרגילים הקודמים,
מתקיים
\[
,e^{2x}\sin x=x+2x^2+\frac{11}{6}x^3+o(x^3)
\]
מכאן
$$.\frac{e^{2x}\sin x+ax+bx^2}{x^3}=\frac{(1+a)x+(2+b)x^2+\frac{11}{6}x^3+o(x^3)}{x^3}$$
ע"מ שהגבול יהיה קיים וסופי נדרוש
$$,a=-1,\ b=-2$$
ואז
\begin{align*}
    \lim_{x\to 0}\frac{e^{2x}\sin x+ax+bx^2}{x^3}=\ \lim_{x\to 0}\frac{\frac{11}{6}x^3+o(x^3)}{x^3}
    =\ \frac{11}{6}.
\end{align*}
\end{solbox}
 \item ידוע שהגבול $\ds{\lim_{x\to 0}\frac{2\sqrt{1+x^2}-x^2-2}{x^\ga}}$ קיים, סופי ושונה מ-$0$. מצאו את $\ga$ ואת הגבול.
\hint{פתחו את $\sqrt{1+t}$ עם $t=x^2$, ומצאו את החזקה הנמוכה ביותר של $x$ שמקדמה במונה אינו מתאפס.}
\begin{solbox}
נזכר כי באחד התרגילים הקודמים קיבלנו כי
$$,\sqrt{1+x^2}=1+\frac{x^2}{2}-\frac{x^4}{8}+o(x^5)$$
מכאן
$$.2\sqrt{1+x^2}-x^2-2=-\frac{x^4}{4}+o(x^5)$$
מתקיים
$$.\lim_{x\to 0}\frac{2\sqrt{1+x^2}-x^2-2}{x^\ga}=\lim_{x\to 0}\frac{-\frac{x^4}{4}+o(x^5)}{x^\alpha}$$
הגבול האחרון קיים וסופי אמ"מ
$\alpha\leq 4$
והוא שווה ל-
$0$
אמ"מ
$\alpha<4$,
לכן
הגבול סופי ושונה מ-
$0$
אמ"מ
$\alpha=4$,
במקרה זה נקבל
$$.\lim_{x\to 0}\frac{2\sqrt{1+x^2}-x^2-2}{x^\ga}=-\frac{1}{4}$$
\end{solbox}
\eE
